\paragraph{Fourier observables of the regulated integral (conditional for \(N\ge3\)).}
Keep the setting, Definition~(A) and the contours \(\mathcal R_u\) of the preceding paragraph, and let \(O\) be a product of \(n_{\mathrm{tr}}\) traces, each of a product of factors \(\mathrm e^{\mathrm i\kappa A_\mu}\) with a single index \(\mu\) and a constant \(\kappa\in\mathbb C\). For real \(\kappa\), \(\operatorname{Tr}\mathrm e^{\mathrm i\kappa A_0}\) is the Fourier mode of \(A_0\) with timelike momentum along the time axis that the regulator fixes. On \(\mathcal R_u\) such an \(O\) is bounded by \(N^{n_{\mathrm{tr}}}\mathrm e^{c\|\tilde A\|}\), with \(c\) the sum of the \(|\kappa|\), so for fixed \(\epsilon\) the Stokes argument above still gives \(I_\epsilon[O]\) on \(\mathcal R_u\) for every \(u<2/3\) (exact).

(i) \emph{Complex momenta.} Let \(0<u<2/3\), with \(\kappa=K\mathrm e^{3\pi\mathrm iu/8}\) in each factor with \(\mu=0\) and \(\kappa=K\mathrm e^{-\pi\mathrm iu/8}\) in each factor with \(\mu\ge1\), each \(K\) real. On \(\mathcal R_u\) every factor is then \(\mathrm e^{\mathrm iK\tilde A_\mu}\), which is unitary, so that \(|O|\le N^{n_{\mathrm{tr}}}\) there. The bound along the flat directions without the polynomial factor and dominated convergence give \(\lim_{\epsilon\to0^+}I_\epsilon[O]=I^{(u)}[O]\), the integral on \(\mathcal R_u\) at \(\epsilon=0\), and the ratios converge where the Euclidean partition function does not vanish. The same construction at \(u=1\) gives the Euclidean expectation value of \(O\) with every factor replaced by \(\mathrm e^{\mathrm iK\tilde A_\mu}\).

(ii) \emph{Real momenta.} For real \(\kappa\neq0\) the integrand on the right-hand side of the limit display above, with \(O=\operatorname{Tr}\mathrm e^{\mathrm i\kappa A_0}\), is not integrable at \(N=2\). There \(|\operatorname{Tr}\mathrm e^{\mathrm iz\tilde A_0}|=2|\cos(z\|\tilde A_0\|_{\mathrm{op}})|\ge2|\sinh(\operatorname{Im}z\,\|\tilde A_0\|_{\mathrm{op}})|\), which grows exponentially for \(z=\kappa\mathrm e^{-3\pi\mathrm i/8}\), while the modulus of the Euclidean weight is the integrand of the bound along the flat directions at \(u=1\), so that, since that integral converges exactly for \(k<14\), and by \(\mathrm{SO}(10)\) invariance, it has no moment of degree \(14\) in \(\|\tilde A_0\|_{\mathrm{op}}\). The polynomial result does not extend to these observables in that form.

(iii) \emph{Smearing in the momentum.} For \(\varphi(\kappa)=\mathrm e^{-(\kappa-\kappa_0)^2/\ell^2}\), with \(\kappa_0\) real and \(\ell>0\), one has \(\int_{\mathbb R}\varphi(\kappa)\operatorname{Tr}\mathrm e^{\mathrm i\kappa A_0}\mathrm d\kappa=\sqrt\pi\,\ell\operatorname{Tr}\mathrm e^{\mathrm i\kappa_0A_0-\ell^2A_0^2/4}\), which on \(\mathcal R_u\) is bounded by \(\sqrt\pi\,\ell N\exp\bigl(\kappa_0^2\sin^2(3\pi u/8)/(\ell^2\cos(3\pi u/4))\bigr)\) for \(u<2/3\). Write \(\mathcal J(\kappa)\) for the limit of \(\langle\operatorname{Tr}\mathrm e^{\mathrm i\kappa A_0}\rangle_\epsilon\) given by (i) for \(\kappa\in\mathrm e^{3\pi\mathrm iu/8}\mathbb R\). Dominated convergence on \(\mathcal R_u\), followed by turning the line of integration in \(\kappa\) for fixed \(\tilde A\) onto \(\mathrm e^{3\pi\mathrm iu/8}\mathbb R\), where the trace is bounded by \(N\) and \(\varphi\) still decays like a Gaussian, gives
\[
\lim_{\epsilon\to0^+}\int_{\mathbb R}\varphi(\kappa)\,\langle\operatorname{Tr}\mathrm e^{\mathrm i\kappa A_0}\rangle_\epsilon\,\mathrm d\kappa=\int_{\mathrm e^{\mathrm i\theta}\mathbb R}\varphi(\kappa)\,\mathcal J(\kappa)\,\mathrm d\kappa\qquad(0<\theta<\pi/4),
\]
and the same holds for \(A_j\), \(j\ge1\), with the lines \(\mathrm e^{-\mathrm i\theta}\mathbb R\), \(0<\theta<\pi/12\).

(iv) \emph{Analytic structure at \(N=2\).} At \(N=2\) the Cauchy--Binet bound of the preceding paragraph holds for every phase and \(\operatorname{Re}\mathrm iS_{\mathrm b}=-\sin(\pi u/2)\,S_{\mathrm E}\) holds on \(\mathcal R_u\) for \(0<u<2\), so the integral of (i) on \(\mathcal R_u\) defines \(\mathcal J(\kappa)\) for \(0<\arg\kappa<3\pi/4\), with \(u=8\arg\kappa/(3\pi)\), and \(|\mathcal J(\kappa)|\le2[\sin(4\arg\kappa/3)]^{-27/2}\). This function is holomorphic. The integrand is even in \(\tilde A_0\), so in polar coordinates for \(\tilde A_0\) the cosine may be replaced by one exponential along full radial lines, and the two halves of each line can be turned to phases slightly below and slightly above \(\arg\kappa\) with the phase of the \(\tilde A_i\) fixed; the integrand is then dominated uniformly for \(\kappa\) near a given point. Hence \(F_{\mathrm E}(t)=\mathcal J(\sqrt t\,\mathrm e^{3\pi\mathrm i/8})\) is holomorphic for \(|\arg t|<3\pi/4\), and for \(t>0\) it is the Euclidean characteristic function \(\langle\operatorname{Tr}\mathrm e^{\mathrm i\sqrt t\,\tilde A_0}\rangle_{\mathrm E}\), which by \(\mathrm{SO}(10)\) invariance is the same for every \(\tilde A_\mu\). For \(A_j\) the same steps give \(F_{\mathrm E}(\mathrm e^{\pi\mathrm i/4}\kappa^2)\) for \(-\pi/4<\arg\kappa<0\). Since \(\arg t=\pi/4\) lies inside the sector, the smeared limit for \(A_j\) in (iii) is \(\int_{\mathbb R}\varphi(\kappa)F_{\mathrm E}(\mathrm e^{\pi\mathrm i/4}\kappa^2)\,\mathrm d\kappa\). For \(A_0\) the line integrals in (iii) are integrals of \(F_{\mathrm E}(\mathrm e^{-3\pi\mathrm i/4}\kappa^2)\), and real \(\kappa\) corresponds to the boundary ray \(\arg t=-3\pi/4\). With \(k^\mu k_\mu=\mp\kappa^2\) for the two cases, both read \(F_{\mathrm E}(\mathrm e^{\pi\mathrm i/4}k^\mu k_\mu)\), the timelike case being reached from \(k^\mu k_\mu-\mathrm i0\). At \(N=2\) the reduction of the preceding paragraph, extended by the angular integral over \(\tilde A_0\), gives the following (numerical). On \(\mathcal R_u\) with \(u=0.6\) and \(u=0.8\) the ratios to \(I^{(u)}[1]\) of the integrals of \(\|\tilde A_0\|_{\mathrm{op}}^{2n}\), \(n=1,2,3\), are the Euclidean moments times \(\mathrm e^{3\pi\mathrm in(u-1)/4}\) to within \(5\times10^{-7}\) relative, and \(\mathcal J\) agrees with the polynomial of degree six in \(\kappa\) built from them to the size of the first omitted term for \(|\kappa|\le1\). For \(u=0.8\) two phases of the \(\tilde A_i\) give the same \(\mathcal J\) to \(5\times10^{-11}\). On \(\mathcal R_{0.6}\) the limits in (i) and (iii) are approached linearly in \(\epsilon\), the ratio of the difference to \(\epsilon\) changing by at most \(2\%\) between \(\epsilon=0.03\) and \(\epsilon=0.003\). The smeared limit at \(\epsilon=0\) is the same on \(\mathcal R_{0.5}\) and \(\mathcal R_{0.6}\) to \(3\times10^{-6}\).

For these observables, this order of limits and \(\gamma=0\), the regulated integral therefore gives at finite \(N\), conditionally for \(N\ge3\), the integrals at \(\epsilon=0\) on the contours \(\mathcal R_u\), and at \(N=2\) values of the continuation \(F_{\mathrm E}\) of the Euclidean characteristic function, or line integrals of them, which \(F_{\mathrm E}\) on \(t>0\) determines. The Euclidean value at phases is not available for real timelike \(\kappa\). This is no statement about \(\langle\operatorname{Tr}\mathrm e^{\mathrm i\kappa A_0}\rangle_\epsilon\) pointwise at real \(\kappa\) as \(\epsilon\to0\); about whether \(F_{\mathrm E}\) continues beyond \(\arg t=-3\pi/4\); about momenta off the coordinate axes, which combine matrices with different phases on \(\mathcal R_u\); about holomorphy for \(N\ge3\); about \(N\to\infty\) taken first; or about \(\gamma\neq0\) and Definitions~(B) and~(C). It does not derive a signature, a background, a frame or the fermion content, and Assumption~\ref{ass:lorentzian} is not discharged. Problem~E remains open and PROGRAMME-CONDITIONAL.
